% [原文 p.139]
\chapter{极小模型纲领的二维奇点}

本章讨论对极小模型纲领而言重要的若干类曲面奇点。

§4.1 给出所有 lc 曲面对的近乎完全的分类。这些结果是指向高维理论的重要路标，其中许多在以后要用到。

§4.2 的目的是用方程对 Du Val 奇点进行分类。这类奇点正是二维标准奇点。从这一视角出发的分类在 [DV34] 中给出。Du Val 奇点在许多不同的情形中自然出现。综述可参阅 [Dur79]。

Du Val 奇点的同时奇点消解在 §4.3 中讨论。这一结果由 [Bri66, Bri71] 确立。这里的证明以 [Tyu70] 为基础。

§4.2–§4.3 在局部解析几何的框架下处理最为自然。任何情形下都可以只用代数簇来处理，但需要作一些颇不自然的构造。因此，在应用这些结果时，我们不得不依赖代数几何与解析几何之间的某些基本比较定理。

§4.4 讨论二维椭圆奇点。这一理论经 [Rei76, Lau77] 的贡献而发展。这些结果在 §5.3 处理三维标准奇点时不可或缺。

§4.5 构造孤立超曲面奇点的半普遍形变空间。

% [原文 p.140]
这些结果全都非常有用，对三维极小模型纲领而言目前不可或缺。另一方面，从构建一般理论的立场看，它们却是障碍。把三维问题归结为二维问题的论证，以及这一过程中用到的曲面论证，大多不能推广到高维。因此，若极小模型纲领在一切维数上都成立，其证明就不得不沿着不同的路线推进。

\section{二维对数标准奇点}
\takeaway{本节完整刻画二维对数标准（lc）对 $(X,\Delta)$ 的局部结构。工具是极小消解上的差异系数：借负定性与伴随公式确定各例外除子的系数，从而把 lc、dlt 曲面对分成若干显式类型。}

本节的目标是详细描述对数标准（lc）二维对 $(X,\Delta)$ 的局部结构。此时完整的分类是已知的。我们将处理其中大部分，但会略去一些以后不需要的细微之处。这里的论证在代数情形与解析情形都有效，而且多数情形下无需修改。

\begin{notation}{4.1}
设 $X$ 是正规代数曲面或正规复解析曲面，$\Delta$ 是 $X$ 上的 $\Q$-除子。此时不必假设 $K_X+\Delta$ 是 $\Q$-Cartier，也可如下定义差异系数。

设 $f\colon Y\to X$ 是 $X$ 的奇点消解，$E=\sum E_i$ 是例外除子。由引理 3.40，相交矩阵 $(E_i\cdot E_j)$ 负定，故关于 $a_i$ 的线性方程组
\[
E_j\cdot\Bigl(\sum_i a_iE_i\Bigr)=E_j\cdot\bigl(K_Y+f_*^{-1}\Delta\bigr),\qquad \forall j
\]
有唯一解。这里将其写成
\[
K_Y+f_*^{-1}\Delta\equiv_f\sum_i a_iE_i .
\]
如此确定的数 $a(E_i,X,\Delta):=a_i$ 称为 $E_i$ 关于 $(X,\Delta)$ 的差异系数（discrepancy）。若 $K_X+\Delta$ 是 $\Q$-Cartier，则它与先前的定义 2.22 一致。

若 $\Delta$ 有效，且对任意 $f$ 及其每个例外除子 $E_i$ 都有 $a(E_i,X,\Delta)\ge -1$，则称 $(X,\Delta)$ 是数值 lc（numerically lc）的。

若 $\Delta$ 有效，且可选取有限集合 $Z\subset X$，使得 $X\setminus Z$ 非奇异，$\Delta|_{X\setminus Z}$
% [原文 p.141]
是 snc 除子，并且对映入 $Z$ 的每个例外除子 $E$ 都有 $a(E,X,\Delta)>-1$，则称 $(X,\Delta)$ 是数值 dlt（numerically dlt）的。

后来的讨论会证明：对曲面而言，若 $(X,\Delta)$ 是数值 lc（或数值 dlt）的，则 $K_X+\Delta$ 是 $\Q$-Cartier 的。亦即 $(X,\Delta)$ 是 lc（或 dlt）的（参见命题 4.11）。因此，“数值 lc”等概念只是为暂时的方便而引入的。

与 §2.3 不同，本章始终假设 $\Delta\ge 0$。这是因为在曲面上，$\Delta\ge 0$ 这一条件在极小奇点消解下得以保持，而且对分类等目的而言 $\Delta\ge 0$ 不可或缺。\hfill$\square$
\end{notation}

% [Tutor Note: 数值 lc 绕开 $K_X+\Delta$ 是 Q-Cartier 这一前提：先借负定性在解消上解方程定出各 $a_i$，再证明曲面情形下 Q-Cartier 自动成立（命题 4.11）。高维没有这种便利，故必须假设 Q-Cartier。]

紧接的 2 个结果是引理 3.41 的直接推论，但因频繁引用，故作为推论陈述如下。

\begin{corollary}{4.2}
设 $U$ 是非奇异曲面，$C=\bigcup C_i$ 是 $U$ 上的连通固有曲线，且相交矩阵 $(C_i\cdot C_j)$ 负定。设 $A=\sum a_iC_i$、$B=\sum b_iC_i$ 是曲线 $C_i$ 的 $\Q$-线性组合。若对任意 $j$ 都有 $(B\cdot C_j)\le(A\cdot C_j)$，则或者对一切 $i$ 有 $a_i=b_i$，或者对一切 $i$ 有 $a_i<b_i$。
\end{corollary}

\begin{corollary}{4.3}
记号同 4.1。若 $K_Y+f_*^{-1}\Delta$ 是 $f$-nef 的，则或者对一切 $i$ 有 $a_i=0$，或者对一切 $i$ 有 $a_i<0$。
\end{corollary}

把标准曲面奇点加以分类，得到的恰好是 Du Val 奇点。这类奇点最早在一系列论文 [DV34] 中得到研究。§4.2 将更详细地考察它们。这里先给出如下定义最为方便。

\begin{definition}{4.4}
设 $(0\in X)$ 是正规曲面奇点的芽，$f\colon Y\to X$ 是极小奇点消解。若对 $Y$ 中每条例外曲线 $E_i$ 都有 $K_Y\cdot E_i=0$，则称 $(0\in X)$ 是 Du Val 奇点。（在定理 4.20 中将证明此时 $K_X$ 是 Cartier 的。）
\end{definition}

\begin{theoremp}{4.5}
记号同 4.1。则下列成立。
\begin{enumerate}[label=(\arabic*),leftmargin=*]
\item $(X,\Delta)$ 是终端的，当且仅当 $X$ 非奇异且 $\mult_0\Delta<1$。
\item $(X,\Delta)$ 是标准的，当且仅当
\[
\begin{cases}
X\ \text{非奇异且}\ \mult_0\Delta\le 1,\ \text{或}\\
(0\in X)\ \text{是}\ \text{Du Val}\ \text{奇点且}\ \Delta=0.
\end{cases}
\]
\end{enumerate}
\end{theoremp}

% [原文 p.142]
\begin{proof}
首先处理 $X$ 非奇异的情形。设 $p\colon X'\to X$，$p^{-1}(0)=E\cong\PP^1$ 表示 $0\in X$ 的爆破。若 $(X,\Delta)$ 是标准的，则得到 $\mult_0\Delta\le 1$（若是终端的，则 $<1$）（引理 2.29）。由于对一切 $q\in p^{-1}(0)$ 都有 $\mult_q p_*^{-1}\Delta\le\mult_0\Delta$，故对爆破次数作归纳，反向（$\Leftarrow$）也得证。

其次设 $X$ 奇异。设 $f\colon Y\to X$ 是 $X$ 的极小奇点消解，$E=\bigcup E_i$ 是例外集，则 $E\ne\varnothing$。若 $(X,\Delta)$ 是标准的，则 $a_i\ge 0$。于是由推论 4.3，对一切 $i$ 有 $a_i=0$。（特别地，若 $(X,\Delta)$ 是终端的，则可知 $X$ 不奇异，故证明完成。）由此 $K_Y+f_*^{-1}\Delta$ 是数值 $f$-平凡的。然而 $K_Y$ 与 $f_*^{-1}\Delta$ 都是 $f$-nef 的，故二者都是数值 $f$-平凡的。因此 $\Delta=0$，且 $(0\in X)$ 是 Du Val 奇点。反之，当 $(0\in X)$ 是 Du Val 奇点时 $K_X$ 是 Cartier 的，这将在定理 4.20 中证明。
\end{proof}

\begin{definition}{4.6}
设 $C=\bigcup C_i$ 是非奇异曲面 $U$ 上的一族固有曲线。$C$ 的对偶图（dual graph）$\Gamma$ 定义如下。
\begin{enumerate}[label=(\arabic*),leftmargin=*]
\item $\Gamma$ 的顶点（vertex）是曲线 $C_i$。
\item 顶点 $C_i$ 附有数字标号（label）$b_i=-(C_i\cdot C_i)$。若已知 $(C_i\cdot C_i)$ 的值，或该值不重要，则省略标号。
\item 顶点 $C_i$ 与 $C_j$（$i\ne j$）由 $(C_i\cdot C_j)$ 条边（edge）相连。
\end{enumerate}
若 $U$ 上还给出不必固有的曲线 $D_i$，则把各 $D_i$ 作为无标号的顶点加入，边的连接仍按上述定义 4.6(3) 处理，由此定义扩张对偶图（extended dual graph）$\Gamma^*$。

% [Tutor Note: 英文本此处作“值未知或不重要时省略标号”，与日文本“已知”相反。底本如此，照译；理解时取“已知该值或该值无关，则省去标号”。]

圈（cycle）是指具有顶点 $C_1,\dots,C_n$ 与边 $e_1,\dots,e_n$（$n\ge 2$）的图，其编号使得 $e_i$ 连接 $C_i$ 与 $C_{i+1}$（其中 $n+1$ 指 $1$）。若 $n=2$，则是两个顶点以及连接它们的两条边。

树（tree）是指任何子图都不是圈的连通图。

% [原文 p.143]
树的顶点 $C_i$，若与它相连的边至少有 $3$ 条，则称为分叉点（fork）。从树中去掉某个分叉点后所得部分的连通分支，称为该分叉点的枝（branch）。

没有分叉点的树称为链（chain）。链的长度定义为顶点的个数。

上述定义尚有不足之处，例如没有排除与自身相交的边。但以后实质使用这些概念时，各 $C_i$ 都是 $\PP^1$，$\bigcup D_j$ 非奇异，$C\cup(\bigcup D_j)$ 是 snc 除子，因此不会产生问题。只是为记号上的方便，单独定义 $n=1$ 的圈。

结点三次曲线（nodal cubic）（或尖点三次曲线（cuspidal cubic））是指具有结点（node）（或尖点（cusp））的不可约三次曲线 $\subset\PP^2$。它在同构意义下唯一确定。$n=1$ 时的圈定义为结点三次曲线。
\end{definition}

\begin{theoremp}{4.7}
设 $(0\in X,\Delta)$ 是数值 lc 对，$f\colon Y\to X$ 是极小奇点消解，$\Gamma$ 是 $f$-例外曲线 $E=\bigcup E_i$ 的对偶图。则下列恰好一条成立。

\begin{enumerate}[label=\arabic*.,leftmargin=*]
\item 例外曲线只有一条 $E=E_1$，它与非奇异椭圆曲线或结点三次曲线同构。此外 $\Delta=0$。此时 $(X,\Delta)$ 不是数值 dlt 的。

\end{enumerate}

在其余情形，各 $E_i$ 都是非奇异有理曲线，且 $\sum E_i$ 是 snc 除子。

\begin{enumerate}[label=\arabic*.,leftmargin=*,start=2]

\item $\Gamma$ 是由非奇异有理曲线构成的圈，且 $\Delta=0$。此时 $(X,\Delta)$ 不是数值 dlt 的。

\item $\Delta=0$，且 $\Gamma$ 是下列之一。（例外曲线的条数为 $m\ge 5$。）

\begin{center}
\begin{tikzpicture}[scale=0.85,every node/.style={inner sep=0pt}]
\draw[thin] (0,0)--(0,1);
\draw[thin] (0,0)--(0,-1);
\draw[thin] (0,0)--(1.25,0);
\draw[thin] (3.75,0)--(5,0);
\draw[thin] (5,0)--(5,1);
\draw[thin] (5,0)--(5,-1);
\draw[thin,fill=white] (0,0) circle (1.8pt);
\draw[thin,fill=white] (0,1) circle (1.8pt);
\draw[thin,fill=white] (0,-1) circle (1.8pt);
\draw[thin,fill=white] (5,0) circle (1.8pt);
\draw[thin,fill=white] (5,1) circle (1.8pt);
\draw[thin,fill=white] (5,-1) circle (1.8pt);
\draw[thin,fill=white] (1.25,0) circle (1.8pt);
\draw[thin,fill=white] (3.75,0) circle (1.8pt);
\node at (2.5,0) {$\cdots\cdots$};
\node at (-0.24,1) {$2$};
\node at (-0.24,-1) {$2$};
\node at (5.24,1) {$2$};
\node at (5.24,-1) {$2$};
\end{tikzpicture}
\end{center}

此时 $(X,\Delta)$ 不是数值 dlt 的。

\item $\Gamma$ 是只有一个分叉点的树，且该分叉点有 $3$ 个枝。

\item （循环商奇点（cyclic quotient singularity））$\Gamma$ 是链。这里允许链为空。

\begin{center}
\begin{tikzpicture}[scale=0.85]
\draw[thin] (0,0)--(1.3,0);
\draw[thin] (2.7,0)--(4,0);
\draw[thin,fill=white] (0,0) circle (2pt);
\draw[thin,fill=white] (4,0) circle (2pt);
\node at (2,0) {$\cdots\cdots$};
\end{tikzpicture}
\end{center}
\end{enumerate}
\end{theoremp}

% [原文 p.144]
\begin{proof}
先从回顾开始。曲线 $C$ 的算术亏格（arithmetic genus）$p_a(C)$ 由 $p_a(C)=1-\chi(\cO_C)$ 给出；对不可约且约化的曲线 $C$，有 $p_a(C)=h^1(\cO_C)\ge 0$，并且下列事实是熟知的。
\[
p_a(C)=0\iff C\cong\PP^1,
\]
\[
p_a(C)=1\iff
\begin{cases}
C\ \text{是非奇异椭圆曲线、结点三次曲线}\\
\text{或尖点三次曲线}.
\end{cases}
\]

现在如记号 4.1 中那样设 $A=\sum a_iE_i$，并如下确定 $d_i$。
\[
d_i:=(A\cdot E_i)=\bigl((K_Y+f_*^{-1}\Delta)\cdot E_i\bigr)=b_i-2+2p_a(E_i)+(E_i\cdot f_*^{-1}\Delta).
\]
由于 $f$ 是极小奇点消解，对一切 $i$ 有 $d_i\ge 0$。在下文的论证中，对每个 $i$ 构造线性组合 $B=\sum b_iE_i$，使对一切 $i$ 都有 $(B\cdot E_i)\le d_i$。由推论 4.2，这意味着 $a_j\le b_j$，即 $A\le B$。而且只要在某一个 $j$ 处取等号，就有 $B=A$。

(1)\ 设 $p_a(E_1)\ge 1$。此时令 $B=-E_1$。这时 $(B\cdot E_i)\le 0\le d_i$（$i\ne 1$）与 $(B\cdot E_1)=b_1\le b_1-2+2p_a(E_1)\le d_1$ 显然成立。因此 $A\le B$。若 $(X,\Delta)$ 是 lc 的，则 $-1\le a_1\le b_1=-1$，故 $A=B$。由于 $\Supp A=E_1$，由推论 4.3 得 $E=E_1$。再由 $b_1=d_1$，得到 $p_a(E_1)=1$ 与 $f_*^{-1}\Delta=0$。若 $E_1$ 有尖点，则再作两次爆破可知 $(X,0)$ 不是数值 lc 的。

% [Tutor Note: 这里的字母 $b_i$ 有双重含义：图 $\Gamma$ 的标号 $b_i=-(E_i\cdot E_i)$，以及组合 $B=\sum b_iE_i$ 的系数。上式中 $-1\le a_1\le b_1=-1$ 用的是后者（$B=-E_1$ 的系数），而 $(B\cdot E_1)=b_1$ 与 $b_1=d_1$ 用的是前者。]

以后可以假设所有 $E_i$ 都是非奇异有理曲线。

(2)\ 设 $(E_1\cdot E_2)\ge 2$。此时令 $B=-E_1-E_2$。作与上面相同的计算，得到 $E=E_1\cup E_2$、$f_*^{-1}\Delta=0$、$(E_1\cdot E_2)=2$。若 $E_1$ 与 $E_2$ 在一点相切，则再作两次爆破可知 $(X,0)$ 不是数值 lc 的。这样便得到 (2) 中 $n=2$ 的情形。

以后可以假设 $\Gamma$ 中两个顶点至多由一条边相连。假设 $\Gamma$
% [原文 p.145]
不是树。即它包含圈，也就是如 (2) 中那样的曲线 $E_1,\dots,E_n$ 的集合。此时令 $B=-E_1-\cdots-E_n$。作与上面相同的计算可得 (2)。若 $n=3$ 且各 $E_i$ 交于一点，则再作一次爆破可知 $(X,0)$ 不是数值 lc 的。

以后可以假设 $E$ 是 snc 除子且 $\Gamma$ 是树。

(3)\ 设 $\Gamma$ 含两个以上的分叉点，或含一个有 $4$ 个以上枝的分叉点。则它包含如下子图。

\begin{center}
\begin{tikzpicture}[scale=0.85,every node/.style={inner sep=0pt}]
\draw[thin] (0,0)--(0,1);
\draw[thin] (0,0)--(0,-1);
\draw[thin] (0,0)--(1.25,0);
\draw[thin] (3.75,0)--(5,0);
\draw[thin] (5,0)--(5,1);
\draw[thin] (5,0)--(5,-1);
\draw[thin,fill=white] (0,0) circle (1.8pt);
\draw[thin,fill=white] (1.25,0) circle (1.8pt);
\draw[thin,fill=white] (3.75,0) circle (1.8pt);
\draw[thin,fill=white] (5,0) circle (1.8pt);
\draw[thin,fill=white] (0,1.05)--(0.05,1)--(0,0.95)--(-0.05,1)--cycle;
\draw[thin,fill=white] (0,-1.05)--(0.05,-1)--(0,-0.95)--(-0.05,-1)--cycle;
\draw[thin,fill=white] (5,1.05)--(5.05,1)--(5,0.95)--(4.95,1)--cycle;
\draw[thin,fill=white] (5,-1.05)--(5.05,-1)--(5,-0.95)--(4.95,-1)--cycle;
\node at (2.5,0) {$\cdots\cdots$};
\end{tikzpicture}
\end{center}

图中 $E_1,\dots,E_4$ 是标有 $\diamond$ 的顶点，$E_5,\dots,E_{4+n}$（$n\ge 1$）是标有 $\circ$ 的顶点。这里令
\[
B=-\frac12(E_1+\cdots+E_4)-(E_5+\cdots+E_{4+n}).
\]
则
\[
(B\cdot E_j)
\begin{cases}
\le 0, & j\ge n+5\ \text{的情形},\\
=b_j-2, & n+4\ge i\ge 5\ \text{的情形},\\
=b_j/2-1, & j=1,\dots,4\ \text{的情形}.
\end{cases}
\]
如此可知，$(X,\Delta)$ 是数值 lc 的，当且仅当 $E=E_1\cup\cdots\cup E_{n+4}$、$f_*^{-1}\Delta=0$、$b_j=2$（$j=1,\dots,4$）全部成立。这就得到 (3)。

% [Tutor Note: 上面分段条件原书写作 $n+4\ge i\ge 5$，而索引 $i$ 在此并未定义；按上下文应为 $n+4\ge j\ge 5$。此处照录，读者勿被 $i$ 误导。]

其余可能性已被定理 4.7 的 (4) 与 (5) 穷尽。
\end{proof}

\begin{remark}{4.8}（关于术语）
(1) 型奇点中 $E=E_1$ 非奇异者称为单纯椭圆型（simple elliptic）。而 (2) 型与 (1) 中结点三次曲线的情形（即圈的情形）合称为尖点（cusp）。尖点奇点可出现在某类模簇 [Hir71] 的边界上。边界点历来被称为尖点。

% [原文 p.146]
(3) 型中 $m\ge 6$ 的奇点可证明是尖点按 $2$ 阶群的商。非奇异点按循环群的商穷尽了 (5) 型奇点。按非循环群的商穷尽了 (4) 型对数终端奇点。只有一个分叉点的对数标准奇点中，非对数终端者可由单纯椭圆奇点作商得到。证明可参阅 [Bri68, KSB88, Lam86]。\hfill$\square$
\end{remark}

\begin{remark}{4.9}
定理 4.7 的逆也成立。即下列结果成立。
\begin{enumerate}[label=(\arabic*),leftmargin=*]
\item 设 $Y$ 是非奇异解析曲面，$E\subset Y$ 是不可约固有曲线。设 $E$ 是椭圆曲线或结点平面三次曲线。若 $(E^2)<0$，则 $E$ 可收缩。设收缩态射为 $f\colon(E\subset Y)\to(0\in X)$，则 $X$ 在 $0$ 处有对数标准奇点。此外，芽 $(0\in X)$ 的解析同构型仅由 $E$ 与 $(E^2)$ 决定。
\item 设 $Y$ 是非奇异解析曲面或代数曲面，$E=\sum E_i\subset Y$ 是连通固有曲线。假设各 $E_i$ 是非奇异有理曲线且 $(E_i^2)\le -2$。再假设 $E$ 的对偶图是定理 4.7(2)–(5) 中所得的图之一，且相交矩阵 $(E_i\cdot E_j)$ 负定。此时 $E$ 可收缩。设其收缩态射为 $f\colon(E\subset Y)\to(0\in X)$，则 $X$ 在 $0$ 处有对数标准奇点。此外，除唯一一个例外，芽 $(0\in X)$ 的解析同构型由 $E$ 的对偶图决定。该例外为
\begin{center}
\begin{tikzpicture}[scale=0.8]
\draw[thin] (0,0)--(1,0);
\draw[thin] (0,0)--(-1,0);
\draw[thin] (0,0)--(0,1);
\draw[thin] (0,0)--(0,-1);
\draw[thin,fill=white] (0,0) circle (2pt);
\draw[thin,fill=white] (1,0) circle (2pt);
\draw[thin,fill=white] (-1,0) circle (2pt);
\draw[thin,fill=white] (0,1) circle (2pt);
\draw[thin,fill=white] (0,-1) circle (2pt);
\end{tikzpicture}
\end{center}
此时同构型由对偶图与分叉点的 $4$ 个枝相交的 $4$ 个点（$\in\PP^1$）的交比决定。
\end{enumerate}
这些结果本书并不需要。klt 情形的可收缩性当 $Y$ 为射影曲面时将在下文证明。其他结果的证明可参阅 [Art62, Art66, Lau73]。
\end{remark}

\begin{propositionp}{4.10}
设 $Y$ 是非奇异射影曲面，$E_i\subset Y$ 是不可约曲线，全体构成 snc 除子 $\bigcup E_i$。设相交矩阵 $(E_i\cdot E_j)$ 负定，且存在有理数 $0\ge a_i>-1$，使对一切 $j$ 有
\[
\Bigl(\bigl(K_Y-\sum_i a_iE_i\bigr)\cdot E_j\Bigr)=0 .
\]
此时存在双有理态射 $f\colon Y\to X$，$X$ 是正规射影曲面，且 $f$ 只收缩 $\bigcup E_i$。此外 $X$ 是 $\Q$-分解的。
\end{propositionp}

% [原文 p.147]
\begin{proof}
令 $\Delta_Y:=-\sum_i a_iE_i$，则 $(Y,\Delta_Y)$ 是 klt 的。取 $Y$ 上的丰富除子 $H$，使 $H-K_Y-\Delta_Y$ 也丰富。存在正数 $b_i\in\Q$，使对一切 $j$ 有
\[
\Bigl(\bigl(H+\sum_i b_iE_i\bigr)\cdot E_j\Bigr)=0 ,
\]
而把 $H$ 换为其适当倍数后，可设 $b_i\in\Z_{>0}$。令 $L:=H+\sum b_iE_i$。此时 $L-K_Y-\Delta_Y$ 是大的。我们来证明 $L$ 与 $L-K_Y-\Delta_Y$ 都是 nef 的。

首先，$(L\cdot E_j)=\bigl((L-K_Y-\Delta_Y)\cdot E_j\bigr)=0$（$\forall j$）成立。又由 $\sum b_iE_i$ 的有效性，对其他曲线 $C\subset Y$ 也有 $(C\cdot L)\ge(C\cdot H)>0$，并且
\[
\bigl(C\cdot(L-K_Y-\Delta_Y)\bigr)\ge\bigl(C\cdot(H-K_Y-\Delta_Y)\bigr)>0 ,
\]
于是 $L$ 与 $L-K_Y-\Delta_Y$ 的 nef 性得证。由定理 3.3，当 $m\gg 1$ 时 $|mL|$ 无基点。取由此诱导的态射的 Stein 分解，得到双有理态射 $f\colon Y\to X$，它恰好收缩 $\bigcup E_j$。此外 $X$ 正规，且双有理变换 $f_*H=f_*L$ 是 $\Q$-Cartier 的。由 Zariski 主定理 [Har77, V.5.2]，$f$ 不依赖于 $L$ 的选取。任取曲线 $C$ 都可表为丰富除子之差 $C=H_1-H_2$，故对 $Y$ 上任意曲线 $C$，$f_*C$ 是 $\Q$-Cartier 的。
\end{proof}

由下面的结果可知，“数值 dlt”与“dlt”等价，“数值 lc”与“lc”等价。

\begin{propositionp}{4.11}
\begin{enumerate}[label=(\arabic*),leftmargin=*]
\item 设 $(X,\Delta)$ 是数值 dlt 对。此时 $X$ 上每个 Weil 除子都是 $\Q$-Cartier 的。
\item 若 $(X,\Delta)$ 是数值 lc 的，则它是 lc 的。
\end{enumerate}
\end{propositionp}

\begin{proof}
若任一情形中 $\Delta\ne 0$，则由推论 4.2，$(X,0)$ 也是数值 dlt 的，故可归结为 (1) 中 $\Delta=0$ 的情形。因此，为证命题 4.11 可假设 $\Delta=0$。设 $f\colon Y\to X$ 是极小奇点消解，$\Delta_Y$ 是 $Y$ 上满足 $K_Y+\Delta_Y\equiv_f 0$ 的 $f$-例外 $\Q$-除子。此时由推论 4.3 有 $\Delta_Y\ge 0$。

(1)\ 先对 $X$ 是射影曲面的情形证明。由数值 dlt 的假设有 $\lfloor\Delta_Y\rfloor=0$。由定理 4.7，$\Delta_Y$ 是 snc 除子。$f$ 是 $E$ 的收缩，故由命题 4.10，$X$ 是 $\Q$-分解的。

对 $X$ 不射影的一般代数情形，把 $X$ 紧化，或在命题 4.10 的论证中用定理 3.24 即可。若 $X$ 是解析曲面，则由定理 4.12 与引理 4.13 得到该结果。

% [原文 p.148]
(2)\ 可设 $(X,0)$ 不是数值 dlt 的，即 $\lfloor\Delta_Y\rfloor\ne 0$。由定理 4.7，$\{\Delta_Y\}$ 是 snc 除子。由于 $-\lfloor\Delta_Y\rfloor\equiv_f K_Y+\{\Delta_Y\}$，故任取 $n\in\Z_{>0}$ 使 $n\Delta_Y$ 为除子，就有
\[
R^1f_*\cO_Y\bigl(n(K_Y+\Delta_Y)-\lfloor\Delta_Y\rfloor\bigr)=0
\]
成立（推论 2.68）。（非射影情形用推论 2.68 的类似结果。）由此得到满射
\[
f_*\cO\bigl(n(K_Y+\Delta_Y)\bigr)\longrightarrow f_*\cO_{\lfloor\Delta_Y\rfloor}\bigl(n(K_Y+\Delta_Y)\bigr).
\]
因此，只要证明
\[
n(K_Y+\Delta_Y)|_{\lfloor\Delta_Y\rfloor}\sim 0 ,
\]
就得到 $n(K_Y+\Delta_Y)\sim_f 0$，从而 $nK_X=f_*(n(K_Y+\Delta_Y))$ 是 Cartier 除子。这一条件可由定理 4.7 如下确认。由分类结果，只可能出现两种情形：$\lfloor\Delta_Y\rfloor$ 是圈且 $\Delta_Y=\lfloor\Delta_Y\rfloor$；或者 $\lfloor\Delta_Y\rfloor$ 是由有理曲线构成、对偶图为树的 snc 除子。前者的情形，$K_{\lfloor\Delta_Y\rfloor}\sim 0$ 成立，故 $n=1$ 即可。后者的情形，$H^1(\cO_{\lfloor\Delta_Y\rfloor})=0$，故选取 $n>0$ 使 $n(K_Y+\Delta_Y)$ 为数值平凡的 Cartier 除子，便有 $n(K_Y+\Delta_Y)|_{\lfloor\Delta_Y\rfloor}\sim 0$（引理 4.13）。
\end{proof}

下面的结果是定理 5.22 的一个非常特殊的情形。

\begin{theoremp}{4.12}
设 $(0\in X,\Delta)$ 是数值 dlt 曲面对，$f\colon Y\to X$ 是极小奇点消解。此时 $R^1f_*\cO_Y=0$。
\end{theoremp}

\begin{proof}
由推论 4.2，$(X,\varnothing)$ 也是数值 dlt 的。设 $E=\sum E_i$ 是 $f$-例外曲线。由定理 4.7，$E$ 是 snc 除子。设 $K_Y\equiv\sum a_iE_i$，则 $0\ge a_i>-1$。其中第一个不等号来自推论 4.3，第二个不等号来自数值 dlt 性的假设。注意 $\cO_Y\equiv K_Y+\sum(-a_i)E_i$。若 $X$ 是代数的，则应用推论 2.68 得 $R^1f_*\cO_Y=0$。

% [Tutor Note: $R^1f_*\cO_Y=0$ 即 $X$ 具有有理奇点。本节末尾的 4.10–4.12 把“数值 dlt / lc”翻译成 Q-Cartier 性与上同调消没，第 5 章的一般理论正建立在这些结论之上。]

解析情形下，或证明推论 2.68 的类似结果，或按 [Art66] 的方式推进论证，二者之一即可完成证明。
\end{proof}

\begin{lemmap}{4.13}
设 $U$ 是（不必约化的）解析空间，且 $H^1(U,\cO_U)=0$ 成立。$Z_U$ 表示取值于 $\Z$ 的常层。此时下列结论成立。
\begin{enumerate}[label=(\arabic*),leftmargin=*]
\item 存在单射 $\Pic(U)\hookrightarrow H^2(U,\Z_U)$。
\item 假设存在紧曲线 $C=\bigcup C_i\subset U$，使 $\mathrm{red}(U)$ 以 $C$ 为收缩核（retract）。此时 $U$ 上线丛 $L$ 的同构类仅由数值 $(L\cdot C_i)$ 决定。
\end{enumerate}
\end{lemmap}

% [块边界：引理 4.13 的陈述起于原书 p.148 末行，止于本块；其证明自 p.149 起，归 ch04b。]
